\chapter{与 D-模的关系}
\markboth{第 5 章\ 与 D-模的关系}{第 5 章\ 与 D-模的关系}

本节把单子 $\oblv^r_X\circ\indf^r_X$ 与 $\oblv^l_X\circ\indf^l_X$ 同微分算子层联系起来。作为结果，我们把 DG 概形上的范畴 $\Crys^r$ 与 D-模的（导出）范畴联系起来。

\section{用积分变换表示晶体}

本小节中，设 $X$ 是局部几乎有限型的 DG ind-概形。

\ob{5.1.1}{n:5.1.1}回忆：对 $X\in\DGIndSch_{\mathrm{laft}}$，范畴 $\IndCoh(X)$ 可对偶化，且有典范的自对偶性，见 [GR1, Sect.~2.6]。

于是对 $X,Y\in\DGIndSch$，范畴 $\Funct_{\mathrm{cont}}(\IndCoh(X),\IndCoh(Y))$ 与
\[
\IndCoh(X)^{\vee}\otimes\IndCoh(Y)\simeq\IndCoh(X\times Y)
\]
等同。具体地说，对象 $Q\in\IndCoh(X\times Y)$ 通过
\begin{equation}
F\mapsto (p_2)_*^{\IndCoh}\circ(\Delta_X\times\mathrm{id}_Y)^!(F\boxtimes Q)\tag{5.1}
\end{equation}
定义函子 $F_Q:\IndCoh(X)\to\IndCoh(Y)$，其中 $p_2:X\times Y\to Y$ 是投影映射，$\Delta_X$ 是对角映射 $X\to X\times X$。

特别地，自函子 $\oblv^r_X\circ\indf^r_X$ 定义了一个对象，记为
\[
\mathcal{D}^r_X\in\IndCoh(X\times X).
\]
我们将确定这个对象。

\ob{5.1.2}{n:5.1.2}记 $\widehat{\Delta}_X$ 为映射
\[
X\times_{X_{\dR}}X\simeq(X\times X)^{\wedge}_X\to X\times X.
\]

\begin{othm}{命题}{5.1.3}{n:5.1.3}
在 $\IndCoh(X\times X)$ 中有典范同构
\[
\mathcal{D}^r_X\simeq(\widehat{\Delta}_X)_*^{\IndCoh}\big(\omega_{X\times_{X_{\dR}}X}\big).
\]
\end{othm}

\begin{proof}
我们从下述一般观察开始。

设函子 $F\in\Funct_{\mathrm{cont}}(\IndCoh(X),\IndCoh(Y))$ 由对应给出，即有 DG ind-概形的图表
\[
\begin{tikzcd}[column sep=large, row sep=large]
 & Z \arrow[dl, "q_1"'] \arrow[dr, "q_2"] & \\
 X & & Y
\end{tikzcd}
\]
并设 $F:=(q_2)_*^{\IndCoh}\circ q_1^!$。设
\[
i:Z\to X\times Y
\]
为诱导的乘积映射。

\begin{othm}{引理}{5.1.4}{n:5.1.4}
在上述情形，函子
\[
(q_2)_*^{\IndCoh}\circ q_1^!:\IndCoh(X)\to\IndCoh(Y)
\]
由核 $Q=i_*^{\IndCoh}(\omega_Z)$ 给出。
\end{othm}

\begin{proof}[引理 5.1.4 的证明]
有图表，其内方块是笛卡尔的：
\[
\begin{tikzcd}[column sep=large, row sep=large]
 Z \arrow[r, "q_1\times\mathrm{id}_Z"] \arrow[d, "i"'] & X\times Z \arrow[r] \arrow[d, "\mathrm{id}_X\times i"'] & X \\
 X\times Y \arrow[r, "\Delta_X\times\mathrm{id}_Y"'] & X\times X\times Y \arrow[ur, "p_2"'] &
\end{tikzcd}
\]
其内方块是笛卡尔的。
对 $F\in\IndCoh(X)$，有
\[
(q_2)_*^{\IndCoh}\circ q_1^!(F)\simeq(p_2)_*^{\IndCoh}\circ i_*^{\IndCoh}\circ(q_1\times\mathrm{id}_Z)^!(F\boxtimes\omega_Z).
\]
由 [GR1, Proposition 2.9.2]\footnote{严格地说，[GR1, Proposition 2.9.2] 中的基变换同构只对竖直箭头为 ind-固有的情形陈述，这翻译成 $i$ 固有。在命题 5.1.3 的证明中，我们正是在这样的情形应用它。}，
\[
i_*^{\IndCoh}\circ(q_1\times\mathrm{id}_Z)^!(F\boxtimes\omega_Z)\simeq(\Delta_X\times\mathrm{id}_Y)^!\big(F\boxtimes i_*^{\IndCoh}(\omega_Z)\big).
\]
\end{proof}

我们把这个引理用于命题 5.1.3 的证明：

由命题 3.3.2(b)，函子 $\oblv^r_X\circ\indf^r_X$ 由对应
\[
\begin{tikzcd}[column sep=large, row sep=large]
 & (X\times X)^{\wedge}_X \arrow[dl, "p_s"'] \arrow[dr, "p_t"] & \\
 X & & X
\end{tikzcd}
\]
给出。断言现在由引理 5.1.4 推出。
\end{proof}

\ob{5.1.5}{n:5.1.5}作为命题 5.1.3 的推论，得到：

\begin{othm}{推论}{5.1.6}{n:5.1.6}
存在典范同构 $\sigma(\mathcal{D}^r_X)\simeq\mathcal{D}^r_X$，其中 $\sigma$ 是交换 $X\times X$ 两因子的换位。
\end{othm}

\section{其他函子的显式公式}

本小节中，设 $X$ 是几乎有限型的最终余连接拟紧 DG 概形。

\ob{5.2.1}{n:5.2.1}现在回忆：范畴 $\QCoh(X)$ 也是紧生成且自对偶的。在同构
\[
\QCoh(X)^{\vee}\simeq\QCoh(X),\qquad \IndCoh(X)^{\vee}\simeq\IndCoh(X)
\]
下，函子 $\Upsilon_X$ 的对偶是 [IndCoh, Sect.~1.1.5] 中的函子 $\Psi_X$（对偶性陈述见 [IndCoh, Proposition 9.3.3]）。

特别地，对 $\mathcal{C}'\in\DGCat_{\mathrm{cont}}$，有
\[
\Funct_{\mathrm{cont}}(\QCoh(X),\mathcal{C}')\simeq\QCoh(X)\otimes\mathcal{C}',
\]
其公式与 (5.1) 类似。

\ob{5.2.2}{n:5.2.2}设 $\mathcal{C}$ 是下列范畴之一：
\[
\QCoh(X\times X)\simeq\QCoh(X)\otimes\QCoh(X),\quad \IndCoh(X\times X)\simeq\IndCoh(X)\otimes\IndCoh(X),
\]
\[
\QCoh(X)\otimes\IndCoh(X)\quad\text{或}\quad\IndCoh(X)\otimes\QCoh(X).
\]
则 $\mathcal{C}$ 是 $\QCoh(X\times X)$ 上的模，我们定义 $\mathcal{C}$ 的自函子
\[
F\mapsto F\{X\},
\]
为与对象
\[
\Cone\big(\mathcal{O}_{X\times X}\to j_*j^*(\mathcal{O}_{X\times X})\big)[-1]
\]
作张量积，其中 $j$ 是开浸入 $X\times X-X\hookrightarrow X\times X$。

注意由 [IndCoh, Proposition 4.1.7 and Corollary 4.4.3]，对 $\mathcal{C}=\IndCoh(X\times X)$，这个函子与
\[
(\widehat{\Delta}_X)_*^{\IndCoh}\circ(\widehat{\Delta}_X)^!
\]
等同，其中 $\widehat{\Delta}_X$ 表示映射
\[
X\times_{X_{\dR}}X\simeq(X\times X)^{\wedge}_X\to X\times X.
\]

\ob{5.2.3}{n:5.2.3}我们断言：

\begin{othm}{命题}{5.2.4}{n:5.2.4}
(a) 与 $\oblv^l_X\circ\indf^l_X$ 相对应、属于
\[
\QCoh(X)\otimes\QCoh(X)\simeq\Funct_{\mathrm{cont}}(\QCoh(X),\QCoh(X))
\]
的对象，典范地等同于
\[
(\Psi_X(\omega_X)\boxtimes\mathcal{O}_X)\{X\}.
\]

(b) 与 $\oblv^r_X\circ\Upsilon_{X_{\dR}}\circ\indf^l_X$ 相对应、属于
\[
\QCoh(X)\otimes\IndCoh(X)\simeq\Funct_{\mathrm{cont}}(\QCoh(X),\IndCoh(X))
\]
的对象，典范地等同于
\[
(\Psi_X(\omega_X)\boxtimes\omega_X)\{X\}.
\]

(c) 与 $\oblv^l_X\circ(\Upsilon_{X_{\dR}})^{-1}\circ\indf^r_X$ 相对应、属于
\[
\IndCoh(X)\otimes\QCoh(X)\simeq\Funct_{\mathrm{cont}}(\IndCoh(X),\QCoh(X))
\]
的对象，典范地等同于
\[
(\omega_X\boxtimes\mathcal{O}_X)\{X\}.
\]

(d) 与 $\oblv^r_X\circ{}'\indf^r_X$ 相对应、属于
\[
\QCoh(X)\boxtimes\IndCoh(X)\simeq\Funct_{\mathrm{cont}}(\QCoh(X),\IndCoh(X))
\]
的对象，典范地等同于
\[
(\mathcal{O}_X\boxtimes\omega_X)\{X\}.
\]

(e) 与 ${}'\oblv^r_X\circ\indf^r_X$ 相对应、属于
\[
\IndCoh(X)\boxtimes\QCoh(X)\simeq\Funct_{\mathrm{cont}}(\IndCoh(X),\QCoh(X))
\]
的对象，典范地等同于
\[
(\omega_X\boxtimes\Psi_X(\omega_X))\{X\}.
\]

(f) 与 ${}'\oblv^r_X\circ{}'\indf^r_X$ 相对应、属于
\[
\QCoh(X\times X)\simeq\Funct_{\mathrm{cont}}(\QCoh(X),\QCoh(X))
\]
的对象，典范地等同于
\[
(\mathcal{O}_X\boxtimes\Psi_X(\omega_X))\{X\}.
\]

(g) 与 $\oblv^r_X\circ\Upsilon_{X_{\dR}}\circ\indf^l_X$ 相对应、属于
\[
\QCoh(X\times X)\simeq\Funct_{\mathrm{cont}}(\QCoh(X),\QCoh(X))
\]
的对象，典范地等同于
\[
(\Psi_X(\omega_X)\boxtimes\Psi_X(\omega_X))\{X\}.
\]

(h) 与 $\oblv^l_X\circ(\Upsilon_{X_{\dR}})^{-1}\circ{}'\indf^r_X$ 相对应、属于
\[
\QCoh(X\times X)\simeq\Funct_{\mathrm{cont}}(\QCoh(X),\QCoh(X))
\]
的对象，典范地等同于
\[
(\mathcal{O}_X\boxtimes\mathcal{O}_X)\{X\}.
\]
\end{othm}

\begin{proof}
设 $\mathcal{C},\mathcal{D}$ 是 $\DGCat_{\mathrm{cont}}$ 的对象，$\mathcal{C}$ 可对偶化，于是
\[
\Funct_{\mathrm{cont}}(\mathcal{C},\mathcal{D})\simeq\mathcal{C}^{\vee}\otimes\mathcal{D}.
\]
设 $F:\mathcal{C}_1\to\mathcal{C}$ 与 $G:\mathcal{D}\to\mathcal{D}_1$ 是连续函子。则诱导的函子
\[
\Funct_{\mathrm{cont}}(\mathcal{C},\mathcal{D})\to\Funct_{\mathrm{cont}}(\mathcal{C}_1,\mathcal{D}_1)
\]
由
\[
(F^{\vee}\otimes G):\mathcal{C}^{\vee}\otimes\mathcal{D}\to\mathcal{C}_1^{\vee}\otimes\mathcal{D}_1
\]
给出。

记住这一点，有：

(a) 与 (c) 由命题 5.1.3、引理 3.4.9 以及下述断言结合推出：
\end{proof}

\begin{othm}{引理}{5.2.5}{n:5.2.5}
伴随 $\mathrm{Id}_{\QCoh(X)}\to\Xi^{\vee}_X\circ\Upsilon_X$ 的单位定义同构
\[
\mathcal{O}_X\to\Xi^{\vee}_X(\omega_X).
\]
\end{othm}

(b) 由命题 5.1.3，利用 $\indf^l_X\simeq(\Upsilon_{X_{\dR}})^{-1}\circ\indf^r_X\circ\Upsilon_X$ 推出。

(d) 由命题 5.1.3，利用同构 ${}'\indf^r_X\simeq\indf^r_X\circ\Xi^{\vee}_X$ 与引理 5.2.5 推出。

(e) 由命题 5.1.3 推出。(f) 由 (d) 推出。(g) 由 (b) 推出。

(h) 由 (d) 推出，利用引理 3.4.9 与引理 5.2.5。

\ob{5.2.6}{n:5.2.6}记 $\mathcal{D}^l_X,\mathcal{D}^{l\to r}_X,\mathcal{D}^{r\to l}_X,\mathcal{D}'^r_X$ 为
\[
\QCoh(X)\otimes\QCoh(X)\simeq\QCoh(X\times X)
\]
中分别对应于函子
\[
\oblv^l_X\circ\indf^l_X,\quad {}'\oblv^r_X\circ\Upsilon_{X_{\dR}}\circ\indf^l_X,\quad \oblv^l_X\circ(\Upsilon_{X_{\dR}})^{-1}\circ{}'\indf^r_X,\quad {}'\oblv^r_X\circ{}'\indf^r_X
\]
的对象。

有：

\begin{othm}{命题}{5.2.7}{n:5.2.7}
对象 $\mathcal{D}^r_X,\mathcal{D}^l_X,\mathcal{D}^{l\to r}_X,\mathcal{D}^{r\to l}_X$ 与 $\mathcal{D}'^r_X$ 之间成立下述关系：

(i) $\mathcal{D}^l_X\simeq(\Psi_X\boxtimes\Xi^{\vee}_X)(\mathcal{D}^r_X)\in\QCoh(X\times X)$；

(ii) $(\Psi_X\boxtimes\mathrm{Id}_{\IndCoh(X)})(\mathcal{D}^r_X)\simeq(\mathrm{Id}_{\QCoh(X)}\boxtimes\Upsilon_X)(\mathcal{D}^l_X)\in\QCoh(X)\otimes\IndCoh(X)$；

(iii) $\mathcal{D}'^{l\to r}_X\simeq(\Psi_X\boxtimes\Psi_X)(\mathcal{D}^r_X)\in\QCoh(X\times X)$；

(iii') $\mathcal{D}'^{l\to r}_X\simeq(\mathrm{Id}_{\QCoh(X)}\boxtimes\Psi_X\circ\Upsilon_X)(\mathcal{D}^l_X)\simeq(\mathcal{O}_X\boxtimes\Psi_X(\omega_X))\otimes_{\mathcal{O}_{X\times X}}\mathcal{D}^l_X\in\QCoh(X\times X)$；

(iv) $\mathcal{D}'^{r\to l}_X\simeq(\Xi^{\vee}_X\boxtimes\Xi^{\vee}_X)(\mathcal{D}^r_X)$；

(v) $\mathcal{D}'^r_X\simeq(\Xi^{\vee}_X\boxtimes\Psi_X)(\mathcal{D}^r_X)\in\QCoh(X\times X)$。
\end{othm}

\begin{proof}
(i) 由引理 3.4.9 推出。

(ii) 由函子的（同义反复的）同构
\[
\oblv^r_X\circ\indf^r_X\circ\Upsilon_X\simeq\oblv^r_X\circ\Upsilon_{X_{\dR}}\circ\indf^l_X\simeq\Upsilon_X\circ\oblv^l_X\circ\indf^l_X
\]
推出。

(iii) 是同义反复的。

(iii') 由 (ii) 与 (iii) 结合推出。

(iv) 由引理 3.4.9 推出。

(v) 是同义反复的。
\end{proof}

\section{关于 t-结构的行为}

继续假设 $X$ 是几乎有限型的拟紧 DG 概形。

\ob{5.3.1}{n:5.3.1}我们注意到：

\begin{othm}{引理}{5.3.2}{n:5.3.2}
对象 $\mathcal{D}^r_X$ 下有界，即属于 $\IndCoh(X\times X)^{+}$。
\end{othm}

\begin{proof}
由命题 5.1.3，利用 $\omega_X\in\IndCoh(X)^{+}$ 以及函子
\[
F\mapsto F\{X\},\qquad \IndCoh(X\times X)\to\IndCoh(X\times X)
\]
右 t-正合推出。
\end{proof}

\ob{5.3.3}{n:5.3.3}现在假设 $X$ 最终余连接。我们断言：

\begin{othm}{命题}{5.3.4}{n:5.3.4}
$\QCoh(X\times X)$ 的对象 $\mathcal{D}^l_X,\mathcal{D}'^{l\to r}_X,\mathcal{D}'^{r\to l}_X,\mathcal{D}'^r_X$ 都最终余连接，即属于 $\QCoh(X\times X)^{+}$。
\end{othm}

\begin{proof}
由命题 5.2.4，利用 $\Psi_X(\omega_X),\mathcal{O}_X\in\QCoh(X)^{+}$ 以及函子
\[
F\mapsto F\{X\},\qquad \QCoh(X\times X)\to\QCoh(X\times X)
\]
右 t-正合推出。
\end{proof}

\ob{5.3.5}{n:5.3.5}最后，假设 $X$ 是光滑经典概形。我们断言：

\begin{othm}{命题}{5.3.6}{n:5.3.6}
对象 $\mathcal{D}^l_X\in\QCoh(X\times X)$ 落在 t-结构的心中。
\end{othm}

\begin{proof}
断言是 Zariski 局部的，因此不妨设 $X$ 仿射。只需证明
\[
(p_2)_*(\mathcal{D}^l_X)\in\QCoh(X)
\]
落在 t-结构的心中。有
\[
(p_2)_*(\mathcal{D}^l_X)\simeq\oblv^l_X\circ\indf^l_X(\mathcal{O}_X)\simeq\Xi^{\vee}_X\circ(\oblv^r_X\circ\indf^r_X)\circ\Upsilon_X(\mathcal{O}_X).
\]
现在，函子 $\oblv^r_X\circ\indf^r_X$ t-正合（见命题 4.2.11），函子 $\Upsilon_X$ 是平移次数 $[n]$ 的等价，而 $\Xi^{\vee}_X$ 是 $\Upsilon_X$ 的逆。
\end{proof}

\section{与微分算子层的关系}

本小节中，取 $X$ 为光滑经典概形。我们将把 $\mathcal{D}^l_X$ 与经典微分算子层 $\mathrm{Diff}_X$ 底层的 $\QCoh(X\times X)$ 中的对象等同起来。

\ob{5.4.1}{n:5.4.1}对任意 $Q\in\QCoh(X\times X)^{\heartsuit}$（它在集合论意义下支撑在对角线上）以及 $F_1,F_2\in\QCoh(X)^{\heartsuit}$，一个映射
\[
(p_2)_*(p_1^*(F_1)\otimes Q)\to F_2
\]
的数据，等价于一个映射
\[
Q\to\mathrm{Diff}_X(F_1,F_2)
\]
的数据。

此外，这个对应与 $\QCoh(X\times X)^{\heartsuit}$ 上的单子结构相容，该结构由卷积与微分算子的复合给出。

\ob{5.4.2}{n:5.4.2}取 $Q=\mathcal{D}^l_X$，$F_1=F_2=\mathcal{O}_X$，由单子 $\oblv^l_X\circ\indf^l_X$ 在 $\mathcal{O}_X$ 上的作用，得到所需的映射
\begin{equation}
\mathcal{D}^l_X\to\mathrm{Diff}_X,\tag{5.2}
\end{equation}
它与代数结构相容。

我们断言：

\begin{othm}{引理}{5.4.3}{n:5.4.3}
映射 (5.2) 是代数的同构。
\end{othm}

\begin{proof}
只需证明 (5.2) 在 $\QCoh(X\times X)$ 的底层对象层面是同构。后者例如由命题 5.2.4 给出的 $\mathcal{D}^l_X$ 作为拟凝聚层的描述推出。
\end{proof}

\section{晶体与 D-模的关系}

设 $X$ 是有限型经典概形。我们将证明：范畴 $\Crys^r(X)$ 可以典范地等同于 $X$ 上右 D-模的（导出）范畴 $\mathrm{D\text{-}mod}^r(X)$。

\begin{othm}{注}{5.5.1}{n:5.5.1}
范畴 $\mathrm{D\text{-}mod}^r(X)$ 满足 Zariski 下降。因此，由命题 3.2.2，下文只需对仿射概形建立典范等价。
\end{othm}

\ob{5.5.2}{n:5.5.2}设 $Z$ 是光滑经典仿射概形，$i:X\hookrightarrow Z$ 是闭浸入。由经典 Kashiwara 引理与命题 2.5.6，为构造等价
\[
\Crys^r(X)\simeq\mathrm{D\text{-}mod}^r(X),
\]
只需对 $Z$ 构造。

于是可以假设 $X$ 本身是光滑经典仿射概形。我们将构造所求的等价，连同函子的交换图
\[
\begin{tikzcd}[column sep=large, row sep=large]
 \Crys^r(X) \arrow[r] \arrow[d, "\oblv^r_X"'] & \mathrm{D\text{-}mod}^r(X) \arrow[d] \\
 \IndCoh(X) \arrow[r, "\Psi_X"'] & \QCoh(X),
\end{tikzcd}
\]
其中右竖直箭头是自然的遗忘函子，函子 $\Psi_X$ 是 [IndCoh, Lemma 1.1.6] 中的等价。

由命题 2.4.4，按上述方式构造等价 $\Crys^r(X)\simeq\mathrm{D\text{-}mod}^r(X)$，等同于构造左晶体与左 D-模之间的等价，连同函子的交换图
\begin{equation}
\begin{tikzcd}[column sep=large, row sep=large]
 \Crys^l(X) \arrow[rr] \arrow[dr, "\oblv^l_X"'] & & \mathrm{D\text{-}mod}^l(X) \arrow[dl] \\
 & \QCoh(X) &
\end{tikzcd}\tag{5.3}
\end{equation}

\ob{5.5.3}{n:5.5.3}由命题 4.7.3、2.4.4 与 4.4.2，范畴 $\Crys^l(X)$ 与其 t-结构心的导出范畴等同。范畴 $\mathrm{D\text{-}mod}^l(X)$ 按定义是 $\mathrm{D\text{-}mod}^l(X)^{\heartsuit}$ 的导出范畴。此外，图 (5.3) 中的竖直箭头 t-正合。

因此，只需在相应的 Abel 范畴层面构造所求的等价
\begin{equation}
\begin{tikzcd}[column sep=large, row sep=large]
 \Crys^l(X)^{\heartsuit} \arrow[r] \arrow[d, "\oblv^l_X"'] & \mathrm{D\text{-}mod}^l(X)^{\heartsuit} \arrow[d, "\mathrm{Id}"] \\
 \QCoh(X)^{\heartsuit} \arrow[r] & \QCoh(X)^{\heartsuit}.
\end{tikzcd}\tag{5.4}
\end{equation}

\ob{5.5.4}{n:5.5.4}后者是 Grothendieck 的经典计算：

具体地说，把 $\Crys^l(X)^{\heartsuit}$ 解释为截断单纯对象
\[
(X\times X\times X)^{\wedge}_X\substack{p_{12}\\[-2pt]p_{13}\\[-2pt]p_{23}}\rightrightarrows(X\times X)^{\wedge}_X\substack{p_1\\[-2pt]p_2}\rightrightarrows X
\]
上的拟凝聚层范畴的心。也就是说，具体地，$\Crys^l(X)^{\heartsuit}$ 的对象是拟凝聚层 $F\in\QCoh(X)^{\heartsuit}$ 连同同构
\[
\varphi:p_2^*(F)\xrightarrow{\sim}p_1^*(F),
\]
它在对角线上限制为恒等，并满足余循环条件
\[
p_{13}^*(\varphi)=p_{12}^*(\varphi)\circ p_{23}^*(\varphi).
\]
下面给出建立 (5.4) 中等价的另一种途径。

\ob{5.5.5}{n:5.5.5}Abel 范畴 $\Crys^l(X)^{\heartsuit}$ 与 $\mathrm{D\text{-}mod}^l(X)^{\heartsuit}$ 分别实现为作用在 $\QCoh(X)^{\heartsuit}$ 上的单子 $M_{\Crys^l(X)}$ 与 $M_{\mathrm{D\text{-}mod}(X)}$ 的模。

按定义，$M_{\mathrm{D\text{-}mod}(X)}$ 由微分算子代数 $\mathrm{Diff}_X$ 给出。单子 $M_{\Crys^l(X)}$ 由 $\oblv^l_X\circ\indf^l_X$ 给出。于是所求等价由引理 5.4.3 推出。

\begin{othm}{注}{5.5.6}{n:5.5.6}
由构造可知，等价
\[
\Crys^l(X)\to\mathrm{D\text{-}mod}^l(X)
\]
与光滑经典概形间映射 $f:Y\to X$ 的拉回相容。
\end{othm}

\begin{annote}
这个等价说明晶体确实推广了 D-模：Kashiwara 引理在晶体的定义中是自动的，而在经典 D-模理论中是需要证明的定理；此外晶体的定义不依赖 $X$ 的光滑性，也不依赖 $X$ 的经典性。\hfill【来源：原书 §0.1.1、§2.5、§5.5（A）】
\end{annote}

